\chapter{Pengujian Sistem}
\label{app:pengujian}

Pengujian sistem \emph{chatbot} hukum dilakukan untuk memastikan bahwa sistem yang dikembangkan berfungsi sesuai dengan kebutuhan yang telah ditetapkan serta mampu menghasilkan jawaban hukum yang relevan dan sesuai dengan konteks peraturan yang berlaku. Pendekatan pengujian yang digunakan adalah \emph{user acceptance test} (UAT), yang menguji sistem dari sudut pandang pengguna akhir berdasarkan kesesuaian perilaku sistem terhadap kebutuhan fungsional dan nonfungsional yang telah didefinisikan.

Dalam pelaksanaan pengujian, sistem menggunakan berbagai pertanyaan hukum melalui antarmuka \emph{chatbot}. Keluaran sistem kemudian diperiksa untuk memastikan bahwa jawaban yang dihasilkan sesuai dengan konteks hukum yang relevan, mencantumkan dasar hukum yang tepat, serta tidak menghasilkan informasi yang menyimpang dari ketentuan peraturan perundang-undangan. Pengujian ini juga mencakup verifikasi pemenuhan kebutuhan fungsional dan nonfungsional sebagaimana telah didefinisikan pada spesifikasi kebutuhan sistem.

Eksperimen kualitas jawaban dilakukan dengan memanfaatkan fitur Klinik Hukum pada platform hukumonline. Pertanyaan-pertanyaan hukum yang tersedia pada fitur tersebut digunakan sebagai data uji, kemudian jawaban yang dihasilkan oleh sistem \emph{chatbot} dibandingkan secara \emph{manual} dengan jawaban rujukan yang tersedia. Proses eksperimen ini dilakukan dengan menilai relevansi substansi hukum, kesesuaian pasal yang dirujuk, serta kelengkapan konteks hukum yang disampaikan.

Berdasarkan hasil pengujian yang dilakukan, sistem \emph{chatbot} hukum mampu menghasilkan jawaban yang relevan dan sesuai dengan konteks hukum yang ada, serta memenuhi kebutuhan fungsional dan nonfungsional yang telah ditetapkan. Dengan demikian, sistem dinyatakan dapat diterima oleh pengguna dan layak digunakan sebagai sistem pendukung dalam penelusuran informasi hukum.
